\documentclass{article}
\usepackage[T1]{fontenc}
\usepackage{lmodern}
\usepackage{graphicx}
\usepackage{amsmath,amssymb}
\usepackage[a4paper,margin=2cm]{geometry}
\usepackage[absolute,overlay]{textpos}
\setlength{\TPHorizModule}{1mm}
\setlength{\TPVertModule}{1mm}
\usepackage[indonesian]{babel}
\usepackage{hyperref}
\setlength{\emergencystretch}{2em}
% unit: Halaman 23

\begin{document}

% === Halaman 23 (llm) ===

\section*{1. Key-Recovery Attack}

\begin{itemize}
    \item Tujuan: memperoleh kunci kriptografi yang digunakan untuk mengenkripsi/dekripsi.
    \item Misalkan sebuah sistem menggunakan
    \[ C = E_K(P) \quad , \quad P = \text{plainteks}, \ K = \text{kunci}, \ C = \text{cipherteks} \]
    \item Dalam \textit{key-recovery attack}, penyerang berusaha menemukan K.
    \item Contoh: Misalkan sebuah cipher menggunakan kunci 16 karakter. Penyerang melakukan pencarian kunci secara \textit{brute force}:
    
    \begin{align*}
    K_1 & \to \text{dekripsi } C \to P_1 \\
    K_2 & \to \text{dekripsi } C \to P_2 \\
    K_3 & \to \text{dekripsi } C \to P_3 \\
    & \vdots
    \end{align*}
    sampai ditemukan kunci yang menghasilkan plainteks yang masuk akal.
\end{itemize}

\end{document}
